% Generated draft: requires mathematical and source review.
\documentclass[11pt]{article}
\usepackage[margin=25mm]{geometry}
\usepackage{fontspec}
\setmainfont{Latin Modern Roman}
\usepackage{amsmath,amssymb,mathtools,unicode-math}
\setmathfont{Latin Modern Math}
\usepackage{enumitem,needspace,xurl,hyperref}
\setlength{\parindent}{0pt}
\setlength{\parskip}{5pt}
\setlength{\emergencystretch}{3em}
\newcommand{\N}{\mathbb{N}}
\newcommand{\Z}{\mathbb{Z}}
\newcommand{\Q}{\mathbb{Q}}
\newcommand{\R}{\mathbb{R}}
\newcommand{\C}{\mathbb{C}}
\newcommand{\F}{\mathbb{F}}
\newcommand{\A}{\mathbb{A}}
\newcommand{\PP}{\mathbb{P}}
\newcommand{\E}{\mathbb{E}}
\newcommand{\eps}{\varepsilon}
\newcommand{\dd}{\,\mathrm{d}}
\begin{document}
\section*{UnsolvedMath ID: 701707}
\textbf{Source problem code:} catalogue entry
\subsection{Definitions and mathematical statement}
For a finite group \(G\), define its exponent by
\[
\exp(G)=\min\{n\in\Z_{>0}:g^n=1_G\text{ for every }g\in G\}
=\operatorname{lcm}\{\operatorname{ord}(g):g\in G\}.
\]
A subgroup is called \(2\)-generated if it can be generated by at most two elements; the generators need not be distinct or nonidentity.

The question is whether
\[
\forall\,G\text{ finite group},\quad \exists\,x,y\in G\quad
\exp(\langle x,y\rangle)=\exp(G).
\]
Here the trivial group is included, with exponent \(1\).

The supplied source states that the answer is affirmative for finite soluble groups. It also states that every finite group has a subgroup generated by at most three elements with the same exponent, attributing this result to E. Detomi and A. Lucchini (2013). These background assertions are retained as supplied source claims and have not been independently verified. The local registry labels the question open; this is a supplied status label, not an independently verified assessment.
\subsection{Short English statement}
Does every finite group contain a subgroup generated by at most two elements whose exponent equals that of the whole group? The supplied source reports an affirmative answer for soluble groups and a three-generator result for arbitrary finite groups.
\subsection{Sources}
\begin{itemize}
\item Local UnsolvedMath registry, supplied complete source record: permanent ID \(701707\), exact source code \texttt{catalogue entry}, title \emph{Kourovka Notebook Problem 20.49}. Supplied status: open. Upstream revision: \texttt{d371ee917ac9647fbb06cd6fc529deec6f3dda82}.
\item \emph{The Kourovka Notebook (Unsolved Problems in Group Theory)}, Problem 20.49, proposed by A. Lucchini in 2022. The supplied record identifies the 20th issue (2022) and the 21st issue edition (2026), arXiv:1401.0300v48. The collection is attributed to editors E. I. Khukhro and V. D. Mazurov.
\item E. Detomi and A. Lucchini, \emph{J. London Math. Soc. (2)} \textbf{87}, no. 3 (2013), 689--706. Supplied reference for the assertion that every finite group has a three-generated subgroup with the same exponent; no article title is supplied.
\end{itemize}
\end{document}
